\documentclass[11pt,a4paper]{article}
\usepackage[utf8]{inputenc}
\usepackage[margin=1in]{geometry}
\usepackage{amsmath,amssymb,amsthm}
\usepackage{listings}
\usepackage{xcolor}
\usepackage{hyperref}

\newtheorem{lemma}{Lemma}

\lstset{
  language=C++,
  basicstyle=\ttfamily\small,
  keywordstyle=\color{blue}\bfseries,
  commentstyle=\color{gray},
  stringstyle=\color{red},
  numbers=left,
  numberstyle=\tiny\color{gray},
  breaklines=true,
  frame=single,
  tabsize=4,
  showstringspaces=false
}

\title{IOI 2018 -- Mechanical Doll}
\author{Solution}
\date{}

\begin{document}
\maketitle

\section{Problem Summary}

Build a switching network of binary switches (Y-shaped connectors called ``switches'') to route a ball through a sequence of triggers $A_1, A_2, \ldots, A_m$ in order. Each switch has two outputs ($X$ and $Y$) and alternates between them each time the ball passes. All triggers route back to the root switch. Minimize the number of switches.

\section{Solution}

\subsection{Key Idea: Complete Binary Tree}

A complete binary tree with $2^k \ge m$ leaves can sequence $m$ triggers. Each traversal from root to leaf visits a unique leaf due to alternating switches. The $2^k$ leaves are visited in \emph{bit-reversal order}: the $i$-th traversal (0-indexed) reaches leaf $\text{bitrev}(i, k)$.

\subsection{Construction}

\begin{enumerate}
  \item Find the smallest $k$ with $2^k \ge m$. Set $L = 2^k$.
  \item The first $L - m$ visits hit \emph{dummy} leaves (routed back to the root switch). The remaining $m$ visits hit $A_1, \ldots, A_m$ in order.
  \item For leaf at tree position $p$ (0-indexed): it is visited at time $t = \text{bitrev}(p, k)$. If $t < L - m$, the leaf is dummy (output $= -1$, i.e., root switch). Otherwise, the leaf outputs trigger $A_{t - (L-m)+1}$.
  \item Build the tree bottom-up. If both children of an internal node are dummy, eliminate that node (replace with $-1$). This reduces the switch count.
  \item Assign switch numbers so the root is switch $-1$ (as required by the grader), with switches numbered $-1, -2, \ldots$
\end{enumerate}

\begin{lemma}
A complete binary tree of depth $k$ with alternating switches visits its $2^k$ leaves in bit-reversal order.
\end{lemma}
\begin{proof}
At depth $d$, a switch sends the ball left on odd passes and right on even passes (or vice versa, depending on initial state). The $i$-th traversal follows the path determined by the bits of $i$ read in reverse, landing at leaf $\text{bitrev}(i, k)$.
\end{proof}

\subsection{Optimization}

Eliminating switches where both children are dummy can reduce the count from $2^k - 1$ to as few as $m - 1$ for favorable $m$.

\section{Implementation}

\begin{lstlisting}
#include <bits/stdc++.h>
using namespace std;

void create_circuit(int M, vector<int> A) {
    int m = A.size();
    if (m == 0) {
        // No triggers, no switches needed
        vector<int> C(M + 1, 0);
        answer(C, {}, {});
        return;
    }

    int k = 0;
    while ((1 << k) < m) k++;
    int leaves = 1 << k;

    // Bit-reverse function
    auto bitrev = [&](int x, int bits) -> int {
        int r = 0;
        for (int i = 0; i < bits; i++) {
            r = (r << 1) | (x & 1);
            x >>= 1;
        }
        return r;
    };

    // Compute leaf targets: leafTarget[visit_time] = output
    int numDummy = leaves - m;
    vector<int> leafTarget(leaves);
    for (int t = 0; t < leaves; t++) {
        leafTarget[t] = (t < numDummy) ? -1 : A[t - numDummy];
    }

    // leafOutput[tree_position] = output for leaf at position p
    vector<int> leafOutput(leaves);
    for (int p = 0; p < leaves; p++) {
        leafOutput[p] = leafTarget[bitrev(p, k)];
    }

    // Build tree bottom-up, eliminating dummy-only subtrees
    // Nodes use 1-indexed complete binary tree layout
    // Leaves at positions [leaves, 2*leaves-1]
    vector<int> nodeOut(2 * leaves);
    for (int p = 0; p < leaves; p++)
        nodeOut[leaves + p] = leafOutput[p];

    int switchCount = 0;
    vector<int> X_out, Y_out;

    for (int node = leaves - 1; node >= 1; node--) {
        int leftOut = nodeOut[2 * node];
        int rightOut = nodeOut[2 * node + 1];

        if (leftOut == -1 && rightOut == -1 && node != 1) {
            nodeOut[node] = -1;  // eliminate this switch
        } else {
            switchCount++;
            nodeOut[node] = -switchCount;
            // X = left child (first visit), Y = right child (second visit)
            X_out.push_back(leftOut == -1 ? -switchCount : leftOut);
            Y_out.push_back(rightOut == -1 ? -switchCount : rightOut);
        }
    }

    // Renumber so root = -1 (currently root = -switchCount)
    auto remap = [&](int x) -> int {
        if (x >= 0 || x < -switchCount) return x;
        return -(switchCount - (-x) + 1);
    };

    for (int i = 0; i < switchCount; i++) {
        X_out[i] = remap(X_out[i]);
        Y_out[i] = remap(Y_out[i]);
    }
    reverse(X_out.begin(), X_out.end());
    reverse(Y_out.begin(), Y_out.end());

    // C[0] = entry point = root switch = -1
    // C[i] for i >= 1: all triggers route back to root
    vector<int> C(M + 1, -1);

    answer(C, X_out, Y_out);
}
\end{lstlisting}

\section{Complexity Analysis}

\begin{itemize}
  \item \textbf{Switches}: At most $2m - 1$ (since $2^k < 2m$). With dummy elimination, often fewer.
  \item \textbf{Time}: $O(m)$.
  \item \textbf{Space}: $O(m)$.
\end{itemize}

\end{document}
